\documentclass[10pt,a4paper]{amsart}
\usepackage[margin=19mm]{geometry}
\usepackage{amsmath,amssymb,mathtools}
\usepackage[scr=rsfs,cal=euler]{mathalfa}
\usepackage{xfrac}
\usepackage[unicode,hidelinks,bookmarks=false]{hyperref}
\newcommand{\ie}{i.e.\ }
\newcommand{\cf}{cf.\ }
\newcommand{\resp}{respectively\ }
\newcommand{\Subsection}{\subsection}
\providecommand{\nobreakdash}{\nobreak\hbox{-}}
\setlength{\emergencystretch}{2em}
\multlinegap0pt
\usepackage{bm}
\usepackage{bbm}
\usepackage{dsfont}
\usepackage{xspace}
\usepackage{cancel}
\usepackage{mathtools}
\usepackage{amsthm}
\makeatletter
\@ifundefined{thm}{\newtheorem{thm}{Thm}}{}
\@ifundefined{prop}{\newtheorem{prop}[thm]{Proposition}}{}
\makeatother
\def\dive{\operatorname{div}}
\title[A fully diagonalized spectral method on the unit ball]{A fully diagonalized spectral method on the unit ball}
\author{Miguel A. Pi\~nar}
\date{Source: arXiv:2601.15911v1 (CC BY 4.0)}
\begin{document}
\maketitle

\noindent\emph{Source and licence.} A fully diagonalized spectral method on the unit ball. Version arXiv:2601.15911v1 (CC BY 4.0), distributed under the Creative Commons Attribution 4.0 International licence, as recorded in the arXiv licence metadata for this exact version and shown on the abstract page https://arxiv.org/abs/2601.15911v1. Full rights notice: licenses/arxiv-2601.15911-CC-BY-4.0.txt.\par\smallskip

\noindent\textit{Editorial scope.} Counted unit: the Prop whose source label is \texttt{prop-3.1} in the article (source0.tex lines 296--307, source label \texttt{prop-3.1}), delivered together with the article's own complete original proof of it (source0.tex lines 310--386, one \texttt{proof} environment of 77 lines). No mathematics is written, repaired or re-derived: statement and proof are the article's own text line for line. Required context retained uncounted: Sob\_ip (lines 105--108). Every definition, display and label of those lines is retained unchanged. Omitted from this package and not redistributed: 1--104, 109--295, 308--309, 387--931. Nothing inside a retained excerpt is omitted or altered: every retained body is delivered exactly as the article's lines are written, so no omission receipt of any kind accompanies this package. Citations: the retained excerpts carry no citation command at all, so no bibliography entry of the article is transcribed and no reference is redistributed. Reference adaptation: none was needed and none was made; every reference inside the delivered unit resolves inside it, so no unresolved reference or citation is produced. Printed slips kept literal: no printed word, symbol, hypothesis or number of the counted statement or of its proof was altered. Layout and class: the article's own class is \texttt{amsart} with packages including amsmath, amsthm, amsfonts, amssymb, graphicx, multirow; the wrapper is the lane's pinned \texttt{amsart} editorial wrapper at 10pt on A4, which restates the theorem-like environments the retained text needs and adds the article's own macro definitions that the retained text needs (\texttt{\textbackslash dive}), with extra wrapper packages (bm, bbm, dsfont, xspace, cancel, mathtools). The delivered numbering is the unit's own. Assets: no retained span contains a figure environment, a graphics-inclusion command, a PDF-page inclusion, a table or a TikZ picture; no image file is copied into this package, the package declares no asset dependency and no image file is redistributed with it. Licence: version arXiv:2601.15911v1 of this article is distributed under the Creative Commons Attribution 4.0 International licence, as recorded in the arXiv licence metadata for this exact version and shown on the abstract page https://arxiv.org/abs/2601.15911v1. A full rights notice is delivered at licenses/arxiv-2601.15911-CC-BY-4.0.txt. Characters: every retained excerpt is pure ASCII, so the delivered engine, XeLaTeX with its default Latin Modern text font, reports no missing character.
\begin{align} 
\langle  u, v\rangle_{S} &:= \mathcal{A} \left((1-\|x\|^2) u(x), (1-\|x\|^2)v(x)\right) =  \lambda\, \int_{{\mathbb B}^{d}} (1-\|x\|^2)^{\kappa+2} u(x) v(x) dx \nonumber  \\ 
  &+ \int_{{\mathbb B}^{d}} \nabla ((1-\|x\|^2)u(x))\,\cdot \nabla ((1-\|x\|^2)v(x)) dx. \label{Sob_ip}
\end{align}

\begin{prop} \label{prop-3.1}
Let $\{Y_\nu^{n-2j}: 1 \le \nu \le a_{n-2j}^d\}$ be an orthonormal basis of $\mathcal{H}_{n-2j}^d$ and let
\[
R_{j,\nu}^{n}(x) := q_{j}(2 \|x\|^2 -1) Y_\nu^{n-2j}(x)
\]
be a basis of Sobolev mutually orthogonal polynomials then $q_{j} = q_{j}^{(\beta)}$ is orthogonal with respect to the univariate \textbf{Sobolev inner product} 
\begin{align}
    (f,g)_{\beta} := & \frac{\lambda}{2^{\kappa+3}} \int_{-1}^1 q_{j}(t) q_{k}(t) (1-t)^{\kappa + 2} (1+t)^{\beta} dt \nonumber\\
 & +    \int_{-1}^1 \frac{d}{dt}\left((1-t)q_j(t)\right) \frac{d}{dt}\left((1-t)q_k(t)\right) (1+t)^{\beta+1} dt, \label{sob_ip_beta} 
\end{align}
where \(\beta = n -2j + \frac{d-2}{2}\). 
\end{prop}

\begin{proof}
From \eqref{Sob_ip} we get
\begin{align}
\langle R_{j,\nu}^{n}(x),R_{k,\eta}^{m}(x)\rangle_{S} =& \lambda\, \int_{\mathbb{B}^d} R_{j,\nu}^{n}(x) R_{k,\eta}^{m}(x) (1-\|x\|^2)^{\kappa + 2} dx \nonumber \\
   &+ \int_{\mathbb{B}^d} \nabla ((1-\|x\|^2)R_{j,\nu}^{n}(x))\,\cdot \nabla ((1-\|x\|^2)R_{k,\eta}^{m}(x)) dx. \label{Sob_ip_mod}
\end{align}
The proof uses the following identity: 
\begin{equation}\label{changevar}
        \int_{\mathbb{B}^d} f(x) dx = \int_0^1 r^{d-1}\int_{\mathbb{S}^{d-1}}f(r\,\xi)\, d\sigma (\xi)\,dr
\end{equation}
that arises from the spherical-polar coordinates $x=r\,\xi$, $\xi\in \mathbb{S}^{d-1}$. In this way, the first integral in \eqref{Sob_ip_mod} gives
\begin{align*}
I_1 =& \int_{\mathbb{B}^d} R_{j,\nu}^{n}(x) R_{k,\eta}^{m}(x) (1-\|x\|^2)^{\kappa+2} dx =\\
   =& \int_{0}^1 q_{j}(2 r^2 -1) q_{k}(2 r^2 -1) (1-r^2)^{\kappa + 2} r^{n-2j+m-2k+d-1} dr \times \\
   &\times \int_{\mathbb{S}^{d-1}}  Y_\nu^{n-2j}(\xi) Y_\nu^{m-2k}(\xi) d\sigma(\xi) \quad \text{ (with } n = m, j = k) \\
=& \int_{0}^1 q_{j}(2 r^2 -1) q_{k}(2 r^2 -1) (1-r^2)^{\kappa + 2} r^{2n-4j+d-1} dr\\ 
& \times \omega_d \delta_{n-2j,m-2k} \delta_{\nu,\eta} \\
=& \int_{0}^1 q_{j}(t) q_{k}(t) \left(\frac{1-t}{2}\right)^{\kappa + 2} \left(\frac{1+t}{2}\right)^{\frac{2n-4j+d-2}{2}}\frac{1}{4} dt\\ 
& \times \omega_d \delta_{n-2j,m-2k} \delta_{\nu,\eta} \\
=& \frac{1}{2^{\kappa+3+n-2j+d/2}} \int_{-1}^1 q_{j}(t) q_{k}(t) (1-t)^{\kappa + 2} (1+t)^{n-2j+\frac{d}{2}-1} dt\\
&\times \omega_d\delta_{n-2j,m-2k} \delta_{\nu,\eta},
\end{align*}
after the change of variable \(t = 2r^2 - 1\).

For the second integral, integration by parts gives
\begin{align*}
I_2 =& \int_{\mathbb{B}^d} \nabla ((1-\|x\|^2)R_{j,\nu}^{n}(x))\,\cdot \nabla ((1-\|x\|^2)R_{k,\eta}^{m}(x)) dx \\
=& - \int_{\mathbb{B}^d} \Delta ((1-\|x\|^2)R_{j,\nu}^{n}(x))(1-\|x\|^2)R_{k,\eta}^{m}(x) dx \\
& + \int_{\mathbb{B}^d} \dive \left(\nabla((1-\|x\|^2)R_{j,\nu}^{n}(x))(1-\|x\|^2)R_{k,\eta}^{m}(x)\right) dx =\\
=& - \int_{\mathbb{B}^d} \Delta ((1-\|x\|^2)R_{j,\nu}^{n}(x))(1-\|x\|^2)R_{k,\eta}^{m}(x) dx,
\end{align*}
which follows from the divergence theorem and the fact that \(1-\|x\|^2\) vanishes on the sphere.

Now, let us write
\begin{align*}
f_j(\|x\|^2) &:= (1-\|x\|^2)  q_{j}(2 r^2 -1)  
\end{align*}
in this way 
\begin{align*}
\Delta &((1-\|x\|^2)R_{j,\nu}^{n}(x)) = \Delta \left( f_j(\|x\|^2) Y_\nu^{n-2j}(x)\right) = \sum_{i=1}^d \partial_i^2 \left( f_j(\|x\|^2) Y_\nu^{n-2j}(x)\right)\\
&= \Delta (f_j(\|x\|^2)) Y_\nu^{n-2j}(x)+ 4 f_j'(\|x\|^2) \langle  x,\nabla\rangle Y_\nu^{n-2j}(x)
+ f_j(\|x\|^2) \Delta Y_\nu^{n-2j}(x) \\
& = \left[ 4 \|x\|^2 f_j''(\|x\|^2)) + 2(2n-4j + d) f_j'(\|x\|^2)\right] Y_\nu^{n-2j}(x),
\end{align*}
since  
\begin{align*}\langle x,\nabla\rangle Y_\nu^{n-2j}(x) &= (n-2j) Y_\nu^{n-2j}(x),\\
\Delta Y_\nu^{n-2j}(x) &= 0,\\
\Delta (f_j(\|x\|^2)) &= 2 d f_j'(\|x\|^2) + 4 \|x\|^2 f_j''(\|x\|^2).
\end{align*}

Thus, using polar coordinates $ x = r \xi$, for the second integral we obtain 
\begin{align*}
I_2 =& - \int_{\mathbb{B}^d} \left[ 4 \|x\|^2 f_j''(\|x\|^2)) + 2(2n-4j + d) f_j'(\|x\|^2)\right] f_k(\|x\|^2) Y_\nu^{n-2j}(x) Y_\eta^{m-2k}(x)dx\\
=& - \int_0^1 \left[ 4 r^2 f_j''(r^2)) + 2(n-2j + d) f_j'(r^2)\right] f_k(r^2) r^{2n-4j+d-1} dr \\
&\times \omega_d \delta_{n-2j,m-2k} \delta_{\nu,\eta}.
\end{align*}
Next, since \(f_k(r^2) r^{2n-4j+d}\) vanishes at the ends of the interval \([0,1]\), integration by parts gives
\begin{align*}
&\int_0^1 4 r f_j''(r^2) f_k(r^2) r^{2n-4j+d} dr =
- 2 \int_0^1  f_j'(r^2) \frac{d}{dr}\left(f_k(r^2) r^{2n-4j+d}\right) dr \\
&= - 4 \int_0^1  f_j'(r^2) f_k'(r^2) r^{2n-4j+d+1} dr  -
2(2n-4j+d) \int_0^1  f_j'(r^2) f_k(r^2) r^{2n-4j+d-1} dr,
\end{align*}
and we get
\[
I_2 = 4 \int_0^1  f_j'(r^2) f_k'(r^2) r^{2n-4j+d+1} dr \omega_d \delta_{n-2j,m-2k} \delta_{\nu,\eta}
\]
Finally, using the change of variables \( t = 2 r^2 -1\), we deduce
\begin{align*}
 4 \int_0^1  f_j'(r^2)& f_k'(r^2) r^{2n-4j+d+1} dr \\
=& \int_0^1 \frac{d}{dr}f_j(r^2) \frac{d}{dr}f_k(r^2)r^{2n-4j+d-1} dr \\
=& \frac{1}{2^{n-2j+d/2}} \int_{-1}^1 \left(- q_j(t) +  (1-t)q_j'(t)\right) \\
& \times \left(- q_k(t) +  (1-t)q_k'(t)\right) (1+t)^{n-2j+d/2} dt \\
=& \frac{1}{2^{n-2j+d/2}} \int_{-1}^1 \frac{d}{dt}\left((1-t)q_j(t)\right) \frac{d}{dt}\left((1-t)q_k(t)\right) (1+t)^{n-2j+d/2} dt,
\end{align*}
and the result follows.
\end{proof}

\end{document}
